\begin{table}[H]
\centering
\caption{Validity of automatically generated routing labels against per-question empirical quality oracles.}
\label{tab:label-validity}
\small
\begin{adjustbox}{max width=\textwidth}
\begin{tabular}{@{}llccccc@{}}
\toprule
Dataset & LLM & Number of Questions & Original Label Agreement (\%) & Tie-Aware Agreement (\%) & Mean F1 Regret & Tie Rate (\%) \\
\midrule
\NA & LLM 1 & \NA & \NA & \NA & \NA & \NA \\
\NA & LLM 2 & \NA & \NA & \NA & \NA & \NA \\
\bottomrule
\end{tabular}
\end{adjustbox}
\TableNote{A/B/C denote No Retrieval, Single-step Retrieval, and Multi-step Retrieval. For $F_{i,s}$ denoting per-question F1, $\mathcal{O}_i=\arg\max_s F_{i,s}$; tie-aware agreement is $\mathbb{1}[y_i\in\mathcal{O}_i]$, and regret is $\max_sF_{i,s}-F_{i,y_i}$. Tie rate is the percentage with $|\mathcal{O}_i|>1$. The cost-aware oracle selects the least expensive quality-equivalent strategy.}
\end{table}
